\section{El mercado chino: velocidad, tecnología y retorno global}

La sección anterior explicó que Mercedes-Benz enfrenta una transformación con múltiples variables; esta sección examina primero el mercado chino. Para Ola Källenius (康林松), las dos palabras clave más importantes sobre el mercado chino son velocidad y tecnología. Los consumidores chinos tienen exigencias muy altas en conducción inteligente, cabina inteligente, interacción por voz y experiencia digital; además, los fabricantes de automóviles y las empresas tecnológicas locales de China iteran con gran rapidez. Para Mercedes-Benz, China ya no es solo un mercado donde se ejecutan localmente las decisiones de la sede alemana: cada vez más innovaciones comienzan en China y después se aprovechan a escala global.

\begin{figure}[H]
\centering
\includegraphics[width=0.96\textwidth]{china-market-flywheel.png}
\caption{Volante de inercia del mercado chino: velocidad china, investigación y desarrollo en China y resultados chinos que retornan al resto del mundo. Diagrama conceptual propio, elaborado a partir de la conversación entre 00:04:12 y 00:18:24.}
\end{figure}

\begin{knowledgebox}{Lectura de la figura: el mercado chino pasa de ser un punto de ejecución a un punto de innovación}
La figura avanza de China speed a Local R\&D, Tech partners, Product fit y Global reuse. Antes, Mercedes-Benz funcionaba más bajo el esquema «Alemania decide, China ejecuta»; ahora Källenius enfatiza que los sistemas de estacionamiento inteligente y de entretenimiento para los asientos traseros desarrollados por el equipo en China pueden llevarse a todo el mundo.
\end{knowledgebox}

\subsection{Qué significa la velocidad china}

Esta subsección descompone la «velocidad china» en varios niveles. El primero es la velocidad de la demanda: los consumidores aceptan con rapidez los vehículos de nuevas energías, la cabina inteligente, la interacción por voz y la integración del ecosistema del teléfono móvil en el automóvil. El segundo es la velocidad de la competencia: las nuevas marcas emergentes y las marcas locales iteran con rapidez, y el precio, el equipamiento, la experiencia y el marketing cambian con alta frecuencia. El tercero es la velocidad de investigación y desarrollo: Mercedes-Benz ya no define sus productos solo en la sede alemana, sino que amplía su capacidad de investigación y desarrollo en Beijing y Shanghái para aprovechar los clústeres tecnológicos de China. El cuarto es el retorno global: las funciones desarrolladas localmente se llevan al resto del mundo.

\noindent\begin{tabularx}{\textwidth}{p{0.18\textwidth}p{0.34\textwidth}X}
\toprule
Dimensión de la velocidad & Manifestación & Exigencia para Mercedes-Benz \\
\midrule
Velocidad de la demanda & Los usuarios aceptan con rapidez la cabina inteligente, la conducción asistida y la experiencia digital & Entender más rápido al cliente chino, no solo traducir y alargar la distancia entre ejes. \\
Velocidad de la competencia & Las nuevas marcas lanzan con alta frecuencia innovaciones de equipamiento y experiencia & No puede iterar solo al ritmo tradicional de los automóviles de lujo. \\
Velocidad de investigación y desarrollo & Los centros de investigación y desarrollo de Beijing y Shanghái asumen más tareas & Los equipos locales necesitan mayor autonomía. \\
Retorno global & El estacionamiento inteligente y el entretenimiento para los asientos traseros desarrollados en China se llevan a todo el mundo & China no es solo un mercado local, también es una fuente de innovación. \\
\bottomrule
\end{tabularx}

\subsection{La batalla de los vehículos eléctricos en China: ¿perdió Mercedes-Benz?}

Esta subsección analiza una pregunta incisiva: la tasa de penetración de los vehículos de nuevas energías en China pasó de 4.7\% en 2019 a casi la mitad en 2024, y Mercedes-Benz enfrenta presión en ventas y participación de mercado de vehículos eléctricos; ¿significa esto que perdió? La respuesta de Källenius es negativa. Subraya que en China la electrificación creció rápidamente a partir de los automóviles urbanos de entrada y de los vehículos de gama media, mientras que la fortaleza de Mercedes-Benz sigue en el segmento de lujo de gama alta; Mercedes-Benz debe ampliar su oferta de vehículos eléctricos, pero también proteger el valor de la marca y el valor residual de los vehículos que ya están en circulación.

\begin{warningbox}{El éxito de una empresa de automóviles de lujo no puede juzgarse solo por la tasa de penetración total}
La tasa de penetración total de los vehículos de nuevas energías en China es muy alta, pero los distintos rangos de precio, usos, grupos de clientes y expectativas de marca no avanzan al mismo ritmo. La cuestión para Mercedes-Benz no es «si electrificarse o no», sino cómo encontrar el ritmo adecuado entre el mercado de gama alta, la conservación del valor de la marca y una oferta completa de vehículos eléctricos.
\end{warningbox}

\begin{importantbox}{Experiencia práctica: el valor para los clientes existentes también es una restricción estratégica}
Källenius menciona repetidamente el valor residual y la conservación del valor. Para una marca de automóviles de lujo, reducir hoy los precios de forma agresiva puede generar ventas a corto plazo, pero también puede dañar a los clientes anteriores, a la red de distribuidores y a la confianza en la marca. Esta es una diferencia fundamental entre la estrategia de las nuevas marcas emergentes y la de una marca centenaria.
\end{importantbox}

\subsection{Las tres etapas de la localización en China}

Esta subsección desglosa con más detalle la idea de «China para el mundo». La primera etapa es la adaptación del producto, por ejemplo los modelos con distancia entre ejes extendida, que consisten en modificar productos globales para que se ajusten mejor a los usuarios chinos. La segunda etapa es la investigación y desarrollo local: Mercedes-Benz amplía sus centros de investigación y desarrollo en Beijing y Shanghái y aprovecha el talento tecnológico y los clústeres industriales de China. La tercera etapa es la exportación global: capacidades como el estacionamiento inteligente y el entretenimiento para los asientos traseros, desarrolladas por el equipo en China, ya no sirven solo a los clientes chinos, sino que pasan a formar parte de los productos globales.

\noindent\begin{tabularx}{\textwidth}{p{0.18\textwidth}p{0.34\textwidth}X}
\toprule
Etapa & Acciones típicas & Significado estratégico \\
\midrule
Adaptación local & Distancia entre ejes extendida, experiencia en los asientos traseros, servicios digitales locales & Responder a las preferencias de los clientes chinos. \\
Investigación y desarrollo local & Ampliación de los centros de investigación y desarrollo de Beijing y Shanghái, incorporación de talento tecnológico & Pasar del punto de venta al punto de innovación. \\
Exportación global & Los sistemas de estacionamiento inteligente y de entretenimiento para los asientos traseros se llevan a todo el mundo & China se convierte en un nodo de la red tecnológica global de Mercedes-Benz. \\
\bottomrule
\end{tabularx}

\begin{importantbox}{El mercado chino no es simplemente un «mercado difícil»}
Si China se ve solo como un mercado muy competitivo y con fuerte presión de precios, se pasa por alto su valor positivo: obliga a las marcas globales a aumentar su velocidad, su experiencia digital y su capacidad de investigación y desarrollo local. Para Mercedes-Benz, el éxito o el fracaso en el mercado chino influye en el ritmo de su transformación global.
\end{importantbox}

\begin{knowledgebox}{Términos clave: tasa de penetración, conservación del valor y oferta de productos}
\noindent\begin{tabularx}{\textwidth}{p{0.18\textwidth}p{0.36\textwidth}X}
\toprule
Concepto & Significado & Papel en este episodio \\
\midrule
Tasa de penetración de los vehículos de nuevas energías & Proporción de vehículos de nuevas energías en las ventas de automóviles nuevos & Muestra la velocidad de cambio del mercado chino. \\
Valor residual & Valor de reventa y de tenencia a largo plazo del vehículo & Una marca de lujo no puede dañar a sus clientes existentes con reducciones de precio extremas. \\
Oferta de productos & El portfolio que cubre distintos rangos de precio y tipos de vehículo & Determina si Mercedes-Benz puede expandirse de la gama alta a un mercado más amplio. \\
Investigación y desarrollo en China & Capacidad de los equipos locales para desarrollar funciones que retornan al resto del mundo & Muestra que el mercado chino se ha convertido en una fuente de innovación. \\
\bottomrule
\end{tabularx}
\end{knowledgebox}

\subsection{La estrategia de lujo no es contraria a la electrificación}

En 2020, Mercedes-Benz anunció que daría más peso a los modelos de lujo de gama alta y de alto margen, lo cual se prestaba a interpretarse como «el lujo tiene prioridad sobre la electrificación». La explicación de Källenius es que el lujo y la electrificación no se oponen. Mercedes-Benz llevará el objetivo de cero emisiones a toda su oferta de productos, incluidos modelos de gama alta como la Clase G, así como el futuro CLA y más modelos de los segmentos principales; el flujo de efectivo que genera el lujo de gama alta se invertirá en tecnología y en nuevos productos.

\begin{figure}[H]
\centering
\includegraphics[width=0.94\textwidth]{luxury-vs-electrification.png}
\caption{Estrategia de lujo frente a estrategia de electrificación: Källenius subraya que no es necesario elegir entre ambas. Diagrama conceptual propio, elaborado a partir de la conversación entre 00:10:38 y 00:13:08.}
\end{figure}

\begin{knowledgebox}{Lectura de la figura: el lujo aporta el flujo de efectivo y la electrificación aporta el camino hacia el futuro}
A la izquierda, Luxury Strategy incluye marca, experiencia, conservación del valor y márgenes de gama alta; a la derecha, Electrification avanza en toda la oferta de productos. El punto de conexión entre ambas es que la solidez financiera equivale a capacidad de innovación: ganar dinero no se aparta de la transformación tecnológica, sino que la hace sostenible.
\end{knowledgebox}

\subsection{Los tres niveles de la estrategia de automóviles de lujo}

La subsección anterior explicó que el lujo y la electrificación no se oponen; esta subsección explica con más detalle la estrategia de lujo en sí misma. La estrategia de automóviles de lujo no consiste solo en vender automóviles caros. El primer nivel es la psicología del cliente: comprar un Mercedes-Benz significa recompensarse y confirmar cierto nivel de vida. El segundo nivel es la experiencia del producto: la seguridad, la calidad, el confort, el diseño, la sensación de manejo y la experiencia en la cabina forman en conjunto la «sensación Mercedes-Benz». El tercer nivel es el modelo de negocio: los modelos de gama alta aportan márgenes y flujo de efectivo, lo que permite a la empresa invertir miles de millones de euros en el desarrollo de nuevas tecnologías. Si solo se mira el precio, se malinterpreta el lujo; si solo se mira la tecnología, también se malinterpreta a Mercedes-Benz.

\begin{knowledgebox}{Términos clave: estrategia de automóviles de lujo}
\noindent\begin{tabularx}{\textwidth}{p{0.18\textwidth}p{0.36\textwidth}X}
\toprule
Nivel & Significado & Papel durante la transformación \\
\midrule
Nivel psicológico & Recompensarse, confirmar la identidad, aspirar a la marca & Mantener el sobreprecio de la marca y la lealtad de los clientes. \\
Nivel de experiencia & Seguridad, calidad, diseño, confort y sensación de manejo & Seguir diferenciándose cuando la tecnología se democratiza. \\
Nivel de negocio & Márgenes de gama alta y flujo de efectivo & Financiar la electrificación, la IA y el desarrollo de plataformas. \\
\bottomrule
\end{tabularx}
\end{knowledgebox}

\subsection{Resumen de la sección}

El mercado chino es la prueba de estrés de la transformación global de Mercedes-Benz. La velocidad, la experiencia inteligente y la investigación y desarrollo local exigen que Mercedes-Benz sea más rápida, pero una empresa de automóviles de lujo no puede sacrificar la marca ni el valor para los clientes existentes en función de un único indicador de ventas.
